\section{The geometry of a finite fluid instruction}
\label{sv:introduction}

A finite list of symbolic instructions becomes a fluid construction only
after several geometric questions have been answered.  Does an instruction
act on every point near a code?  Can its target overlap another source?
How is the input supplied without changing the requested time symmetry?
Does a particle avoid the detector during the motion between two successive
codes?  These questions distinguish constructions that have the same
integer-time symbolic action.

We study finite lists of affine maps between closed rectangles, together
with the local tape rules that produce those maps. Throughout, a rectangle
is an axis-parallel Cartesian product of intervals in the stated coordinates.
A \emph{geometric
instruction} specifies a source rectangle, a target rectangle, and an
affine bijection between them.  Sources in a list must be mutually
separated, and so must targets; intersections between the two lists are
allowed.  A reciprocal diagonal instruction expands one coordinate by
the factor by which it contracts the other.  A general positive diagonal
instruction may change the rectangle's area.  We realize these operations
on sheets or on boxes of positive thickness using smooth transverse
velocity pulses, with explicit control of their entire trajectories.

\subsection{Main results}
Write $\mathbb T_L^3=(\mathbb R/L\mathbb Z)^3$ for the flat torus.
At a fixed viscosity $\nu>0$ we use
\begin{equation}\label{sv:NS}
 u_t+(u\cdot\nabla)u=-\nabla p+\nu\Delta u+f,
 \qquad \operatorname{div}u=0,\qquad u(0,\cdot)=0.
\end{equation}
Pressure is spatially periodic and normalized to have mean zero.
A classical comparison solution has continuous $u,u_t$, spatial
derivatives of $u$ through order two, and $p,\nabla p$ on every finite
closed time cylinder.  The material flow is defined by
$\partial_tX(t,a)=u(t,X(t,a))$, $X(0,a)=a$.
An effective force prescription evaluates every requested mixed derivative
at computably supplied space and time arguments to any positive rational
error.  The construction also supplies effective global bounds when these
are asserted.  Its finite formula describes all instruction branches;
it is not a recording of an executed computation.

\begin{theorem}[Periodic initialization]\label{sv:periodic-theorem}
Fix a computable viscosity $\nu>0$.  From a deterministic one-tape
machine, with head moves in $\{-1,0,1\}$, and a finite input, one can
effectively construct a smooth force on $\mathbb T_1^3$ such that:
\begin{enumerate}[label=(\roman*),itemsep=2pt]
\item The force is divergence free, has zero spatial mean, and is
one-periodic from time zero.  Equation~\eqref{sv:NS} has a global
smooth solution with pressure zero, unique in the stated classical class.
\item The velocity is one-periodic, vanishes near every integer time,
has zero instantaneous self-advection, and has uniformly bounded kinetic
energy.  Every mixed derivative of force and velocity is globally bounded.
\item For the fixed label $a_*=(1/4,1/2,1/4)$ and fixed open strip
$O=\{1/2<x_1<7/8\}$, the material trajectory enters $O$ at a finite
real time exactly when the machine halts, including an initially halted
machine.
\end{enumerate}
The construction may use one Cartesian velocity component at a time.
For arbitrary fixed real $\nu>0$ the same formulas hold and are
effective relative to that viscosity.
\end{theorem}

Section~\ref{sv:periodic-compiler} proves this theorem after the
common toolkit in Section~\ref{sv:toolkit}. It gives two initialization
mechanisms: an input-dependent chart and a reserved loading branch in
the repeated table. The latter gives the one-component option. Both
incorporate initialization into a schedule that repeats from time zero.
Alternative recorder guards and observers are developed separately in
Section~\ref{sv:recorders}.

The second result allows an instruction to change planar area.  Its
normal action then changes as well, as volume preservation requires.

\begin{theorem}[General diagonal maps on a sheet]\label{sv:sheet-theorem}
Let $P_i,Q_i\subset[2,3]^2$, $1\le i\le N$, be closed rectangles
with rational vertices and positive side lengths.  Suppose the sources are pairwise
positively separated and the targets are pairwise positively separated.
Let $F_i:P_i\to Q_i$ be the center-to-center affine map with any
positive rational diagonal linear part, with $F_i(P_i)=Q_i$.  For each computable $\nu>0$
there is an effective smooth mean-zero solenoidal force on
$\mathbb T_{10}^3$, one-periodic from time zero, whose unique solution
in the classical class has pressure zero and period map
\[
 (y,2)\longmapsto(F_i(y),2)\qquad(y\in P_i).
\]
Its velocity is zero near integer times and has zero self-advection;
all mixed derivatives of force and velocity are bounded.  No condition
is imposed on source--target intersections.  For $N=0$ use zero.
\end{theorem}

The sequential storage construction in Section~\ref{sv:storage} proves
this theorem for arbitrary such rectangle data;
Section~\ref{sv:separator} applies it to a complete tape computation.
Section~\ref{sv:fifteen} gives a simultaneous construction for the
narrow-strip rectangles of a different recorder, using separate heights: contraction,
translation, and expansion take fifteen phases.  Both proofs exhibit the
compensating normal action instead of treating an area-changing planar
map as area preserving.  If the required action is on a whole box of
positive thickness, Section~\ref{sv:parking} provides a separate
reciprocal box-transfer theorem with an all-time sign guard.

\subsection{How the constructions fit together}
A finite local recorder stores enough information to reverse each rule
on its entire partial domain. Two positional coordinates encode the tape
halves. Fixing the finitely many letters read by a rule gives a closed
rectangle, and the inverse ensures that the image rectangles are
separated. The toolkit in Section~\ref{sv:toolkit} makes this passage
explicit.

Each reciprocal rectangle map is then performed at a private height.
A vertical pulse lifts the source rectangle; four centered planar shears
produce the diagonal change of shape; a translation moves it to its
target center; a final vertical pulse returns it to the coding height.
Separating these stages in time makes each active field transverse, so
convection vanishes. Spatial selectors have zero mean and are constant
on the tracked paths. Thus the direct force formula gives the fluid
equation, while the partial shear paths give the observation bounds.
Sections~\ref{sv:reciprocal-applications} and~\ref{sv:routing} adapt
this mechanism to a one-time loader, other clocks and charts, and whole
boxes.

For arbitrary planar area change, the normal direction must also move.
A vertical shear records a horizontal offset in height, a horizontal
shear changes that offset, and a second vertical shear restores the
sheet while rescaling nearby normal offsets. To allow a target to
intersect any source, the general construction first evacuates all
source sheets into a separate storage region. It then delivers them one
at a time. This geometry proves the second theorem before any symbolic
application is chosen; its low-height bound supplies the application's
fixed detector.

\subsection{Ideas and relation to earlier work}
An inverse cannot discard the information that distinguishes predecessors.
Landauer discusses this retention principle and its storage implications
\cite[Section~3]{Landauer1961}.  Reversible Turing machines already
appear in Lecerf's work on undecidability for code isomorphisms
\cite{Lecerf1963}.  Bennett's reversible simulation records each
instruction and verifies that the local rules have disjoint domains and
disjoint ranges \cite[Eqs.~(5)--(11) and Table~1]{Bennett1973}.
Our finite local recorders use these principles; their particular guards
are checked on the full partial domains, including tapes outside the
initialized run.

Moore's generalized shifts connect positional tape coding to dynamical
systems \cite{Moore1990}.  They admit finite affine Cantor-block
representations, and invertible generalized shifts have smooth
realizations \cite[Lemma~0 and Section~6]{Moore1991}.
Our geometric constructions implement the resulting closed-rectangle
maps with explicit transverse pulses, mean-zero forcing, and observation
bounds throughout each operation.  Separation of sources and of images
is the interface between the symbolic inverse and these pulse formulas.

The fluid setting also has substantial predecessors.  Cardona, Miranda,
Peralta-Salas and Presas realize computation by stationary Euler flows
on an adapted Riemannian three-sphere \cite{CMPP}.  Their
Proposition~5.1 constructs an area-preserving realization by contracting,
transporting, and expanding neighborhoods with a relative area correction.
That geometric organization is related to the sheet construction below.
For a fixed Euclidean metric, Cardona, Miranda and Peralta-Salas
construct Turing-complete stationary Euler fields on $\mathbb R^3$;
their simulation uses a noncompact set and the fields have infinite
energy \cite{CMP2023}.  On the standard flat three-torus they robustly
simulate tape-bounded machines, with reachability tested before a
trajectory leaves a prescribed compact region.  Their construction also
gives time-dependent Navier--Stokes realizations of these bounded
computations \cite{CMP2023}.

Dyhr, Gonz\'alez-Prieto, Miranda and Peralta-Salas obtain stationary
unforced viscous computation after a metric deformation, using the
Hodge Laplacian on one-forms as the viscous operator
\cite[Theorem~A and Proposition~3.1]{DGMP}.  Here the flat metric,
the usual flat Laplacian, positive viscosity, and zero initial velocity
are fixed; the finite machine and input specify the external force.
The local shear formulas supply the required geometry directly.

The companion \emph{Finite Instructions and Solenoidal Shear Flows}
\cite{Entry} gives one complete initialized construction and the common
reciprocal-rectangle theorem.  We restate its geometric realization
(Theorem~3.2) in Section~\ref{sv:imports}, and its recorder and
extensions (Lemma~2.1 and Section~5.1) in
Section~\ref{sv:recorder-import}.  The present proofs address changed requirements:
periodicity from zero, alternative full-domain guards, routing through a
common center, positive thickness, and planar area change.  A symbolic
inverse is used to prove separation; it does not replace the geometric
or intermediate-time argument.

At each instant our active velocity moves only in directions on which its
coefficients do not depend.  Consequently convection vanishes and
$f=U_t-\nu\Delta U$ gives solenoidal forcing with pressure zero.
Section~\ref{sv:analytic} proves the relevant uniqueness statement by the
classical difference-energy method associated with Leray
\cite[Section~18]{Leray1934}.  Slower onto clocks preserve the complete
trajectory, including initialization; their exact derivative estimates
are proved in Section~\ref{bcs:clocks-section}.  A clock with finite
accumulated time would not suffice.

Finally Section~\ref{sv:material} writes the same equation in material
labels, and Appendix~\ref{sv:interpreter} constructs a fixed universal
table in the stated tape model.  Thus the force compiler can either take
an arbitrary finite machine table or keep one interpreter fixed and vary
only its input.  Composing Theorem~\ref{sv:periodic-theorem} with any
decision procedure for its fixed particle event would decide halting.
This includes Turing's symbol-printing question: stop when the designated
symbol is printed, and otherwise continue forever, including when the
original machine stops without printing it \cite[Section~8]{Turing1936}.
The constructions use exact real codes and assert no uniform
finite-precision tolerance.
